\documentclass[a4paper,12pt]{article}
\usepackage[utf8]{inputenc}
\usepackage{xeCJK}  % 中文支持
\usepackage{ctex}
\usepackage{amsmath,amsthm,amssymb}  % 数学包
\usepackage{graphicx}  % 插入图片
\usepackage{booktabs}  % 美化表格
\usepackage{hyperref}  % 超链接
\usepackage{enumitem}  % 自定义列表
\usepackage[left=3.5cm,right=3cm,top=2cm,bottom=2cm]{geometry}  % 页面设置
 
\title{读原著，说新知——《反对本本主义》}
\author{2023080904007 吕乐彬}

 
\begin{document}
 
\maketitle
 
\section{简介}
毛泽东于1930年写就的《反对本本主义》，
是一篇跨越时空、历久弥新的马克思主义方法论经典。
我首先精读了原典，剖析“本本主义”的实质与危害，理解“没有调查，没有发言权”这一核心命题的深刻内涵。进而，将这一思想置于新时代的语境下，探讨其对于破除“新本本主义”、应对信息茧房、推动高质量发展以及提升个人认知与实践能力的重大启示。《反对本本主义》所倡导的调查研究与实事求是精神，不仅是中国共产党克难制胜的法宝，更是当代社会每一个组织与个体都应秉持的科学方法论。

 
\section{背景}

《反对本本主义》诞生于中国革命的关键时期，其直接目的是纠正党内盛行的教条主义错误。当时，一些同志习惯于机械地执行共产国际的指示和照搬马克思主义的“本本”，严重脱离了中国革命的实际。毛泽东同志以巨大的理论勇气和政治魄力，发出了“反对本本主义”的呐喊。今天，我们重读这篇著作，不应仅仅视其为一段历史文献，更应将其视为一把能够刺破当下诸多现实问题的思想利剑。它教导我们的，不仅是一种工作方法，更是一种世界观和方法论。

\section{核心要义}

\begin{itemize}
    \item \textbf{“没有调查，没有发言权”}：
这是毛泽东在《反对本本主义》中开宗明义提出的根本性原则，是其唯物论与实践观在工作方法上的集中体现。他尖锐地指出：“你对于某个问题没有调查，就停止你对于某个问题的发言权。”并进一步阐发道：“一切结论产生于调查情况的末尾，而不是在它的先头。”这一论断的核心在于，它确立了一个认识论上的严格程序：对客观实际的周密调查，是形成任何正确认识与观点的唯一源泉，是先于一切讨论、决策和行动的不可逾越的前提。它从根本上批判了那种脱离实际、仅凭主观想象或书本教条就妄下结论的唯心主义倾向，为党的实事求是思想路线奠定了坚实的方法论基础。
    \item \textbf{“本本主义”的危害}:
毛泽东将党内存在的教条主义错误思想形象而深刻地命名为“本本主义”。他精准地剖析了其本质与危害：首先，它颠倒了理论与实践的源流关系。毛泽东强调：“我们说马克思主义是对的，决不是因为马克思这个人是什么‘先哲’，而是因为他的理论，在我们的实践中，在我们的斗争中，证明了是对的。”这明确指出，理论的权威性来自于其被实践检验的正确性，而非其本身是神圣不可侵犯的教条。其次，本本主义在现实中表现为对上级指示的盲目执行。他批判道：“盲目地表面上完全无异议地执行上级的指示，这不是真正在执行上级的指示，这是反对上级指示或者对上级指示怠工的最妙方法。” 这揭示出，脱离具体环境、缺乏调查研究的机械执行，实质上是对上级精神的曲解与破坏，最终会窒息革命事业的活力，导致工作的失败。
    \item \textbf{调查研究是“十月怀胎”，解决问题是“一朝分娩”}：
毛泽东运用“十月怀胎”与“一朝分娩”这一极其生动而科学的比喻，深刻揭示了调查研究与解决问题之间的内在逻辑与辩证关系。他明确指出：“调查就像‘十月怀胎’，而解决问题就像‘一朝分娩’。调查就是解决问题。” 这个比喻蕴含了三层核心要义：第一，调查研究是一个必须投入充分时间与精力的、长期且艰苦的孕育过程，需要耐心与细致，无法一蹴而就。第二，解决问题是前期充分调查研究所必然导致的结果，是水到渠成、瓜熟蒂落的自然阶段。第三，它从根本上定义了调查研究的性质——其本身不是目的，而是内在地包含了解决问题的指向；真正的调查研究，其全过程就是为了最终“解决问题”。这一论断彻底划清了科学的调查研究与形式主义的“走过场”式调研的界限。
\end{itemize}
 

\section{新时代启示}

\begin{itemize}
    \item \textbf{在“信息过载”时代，更要捍卫“田野调查”的价值。}
  我们身处一个被数据、报表和屏幕信息包围的时代。然而，海量的二手信息并不能替代亲身实地的“田野调查”。正如毛泽东所倡导的“开调查会”，要“眼睛向下”，深入基层、深入群众。面对复杂的社情民意，坐在办公室里看报告、刷数据，极易产生“信息茧房”和认知偏差。真正的“新知”，往往蕴藏在街头巷尾、工厂车间、田间地头的鲜活实践中。因此，反对当代的“本本主义”，就是要反对对数字屏幕的盲目依赖，重拾“接地气”的调研传统，从火热的生活一线获取最真实、最生动的第一手材料。
    \item \textbf{克服“路径依赖”与“经验主义”，激发创新活力。}
    过去的成功经验，若被奉为不可逾越的“本本”，就会成为阻碍创新的“路径依赖”。许多企业和组织习惯于沿用过去的成功模式，面对新市场、新技术、新群体时反应迟钝，这正是“本本主义”的变种。毛泽东强调：“共产党的正确而不动摇的斗争策略，决不是少数人坐在房子里能够产生的，它是要在群众的斗争过程中才能产生的。”这启示我们，创新不是闭门造车，而是要打破对过往“本本”的迷信，以开放的心态深入新实践，尊重群众的首创精神，在动态的探索中寻找新方法、新路径。

    \item \textbf{破除“唯上”思维，构建实事求是的组织文化。}
    《反对本本主义》中批评了“盲目地表面上完全无异议地执行上级的指示”的现象。这在今天，表现为一种“唯上不唯实”的官僚主义作风。一些干部习惯于将上级文件作为行动的“唯一本本”，缺乏结合本地、本部门实际进行创造性落实的勇气和智慧。构建实事求是的组织文化，就要鼓励干部在深刻领会上级精神的基础上，敢于并善于根据实际情况作出科学决策。评价工作，不能只看“对标对表”的完成度，更要看解决实际问题的“绩效”和群众的“满意度”。

\end{itemize}
 
\section{结论：永不过时的思想武器}
 重读《反对本本主义》，我们深刻体会到，经典原著的价值在于其能够跨越时空，持续提供分析问题和解决问题的基本立场、观点和方法。“读原著”不是寻章摘句，而是为了“说新知”、开新局。在当前全面建设社会主义现代化国家的征程中，我们面临的挑战前所未有，新情况、新问题层出不穷。唯有坚持“没有调查，没有发言权”的原则，大力弘扬实事求是的精神，坚决反对一切形式的本本主义、教条主义和形式主义，才能不断获得真知灼见，找到破解难题的“金钥匙”，从而在理论与实践的结合中，推动事业不断向前发展。《反对本本主义》这盏思想的明灯，必将持续照亮我们前行的道路。

\end{document}